\documentclass{article}
\usepackage{amsmath}
\usepackage{graphicx}
\graphicspath{ {images/} }
\usepackage[spanish]{babel}
\usepackage{bbm}
\usepackage{amsthm}
\usepackage{authblk}
\usepackage{hyperref}
\usepackage[utf8]{inputenc}
\usepackage{mathrsfs}
\usepackage{amsfonts}
\usepackage{amssymb}
\usepackage{enumerate}
\usepackage[left=3cm,right=2cm,top=2cm,bottom=2cm]{geometry}
\usepackage{epsfig}
 \renewcommand{\baselinestretch}{0.93}
 \thispagestyle{empty}
 \pagestyle{empty}
\setlength{\textheight}{53pc} \setlength{\textwidth} {36pc}
\setlength{\topmargin}{-4mm}
\usepackage{multicol}
\newcommand*{\email}[1]{%
    \normalsize\href{mailto:#1}{#1}\par
    }
\setlength{\columnsep}{1cm}
\newcommand{\N}{\mathbb{N}}
\newcommand{\calM}{\cal{M}}
\newcommand{\calP}{\cal{P}}
\newcommand{\F}{\mathbb{F}}
\newcommand{\R}{\mathbb{R}}
\newcommand{\Z}{\mathbb{Z}}
\newcommand{\Q}{\mathbb{Q}}
\newcommand{\C}{\mathbb{C}}
\newcommand{\bbH}{\mathbb{H}}


\theoremstyle{plain}
\newtheorem{thm}{Teorema}[section]
\newtheorem{prop}[thm]{Proposición}
\newtheorem{lemma}[thm]{Lema}
\newtheorem{cor}[thm]{Corolario}
\theoremstyle{definition}
\newtheorem{rec}[thm]{Recuerdo}
\newtheorem{defin}[thm]{Definición}
\newtheorem{ejem}[thm]{Ejemplo}
\newtheorem{ejer}[thm]{Ejercicio}
\newtheorem{sol}[thm]{Solución}
\newtheorem{rmk}[thm]{Observación}
\author{Marcos Morales}
\affil{Universidad Finis Terrae}
\title{Primera Semana\\}






\begin{document}


\begin{picture}(400, 75)
   \put(-40,15){\includegraphics[width=5cm]{UFT.png}}
\end{picture}


\begin{center}
\huge{D\'ecima Semana}
\end{center}

\begin{center}
\Large{ Marcos Morales, Camila Muñoz}
\end{center}

\begin{center}
	\large{Cálculo Vectorial}
\end{center}
\begin{center}
\setlength{\parindent}{0cm}\large{Clases centradas para capacitar a los estudiantes de ingeniería en Cálculo Vectorial. \\}
\end{center}


\vspace{1cm}

\begin{center}
	\large{Contenidos:}
\end{center}
\vspace{0.1cm}
\begin{center}
\begin{minipage}{0.5\textwidth}
\begin{enumerate}
    \item Teorema de Green
\end{enumerate}
\end{minipage}
\end{center}




\vspace{6cm}




\pagestyle{empty}


\setcounter{page}{1}
\pagestyle{plain}
\setlength{\parindent}{0cm}





\newpage


\section{Definiciones}

\begin{defin}Una curva $C$ parametrizada por $r(t)$ con $a \leq t \leq b$, se dice \textbf{cerrada}, si su punto inicial es igual al punto final, es decir $r(a)=r(b)$, si $F$ es un campo vectorial, entonces la integral de línea sobre \textbf{curva cerrada} $C$, se denota por
\end{defin}

\begin{align*}
    \oint_C F \cdot dr
\end{align*}

\begin{rmk}Una integral de línea sobre una curva cerrada es el mismo concepto que integral de línea normal, solo que se elige una notación distinta debido a la relevancia que tienen las curvas cerradas. \\
\end{rmk}

\begin{ejem}La circunferencia $C$ es una curva cerrada, una parametrización es por ejemplo $r(t)=(\cos(t),\sin(t))$, $0 \leq t \leq 2\pi$ \\
\end{ejem}


\begin{defin}Una curva $C$ se dice que está \textbf{orientada positivamente} si su parametrización $r(t)$, recorre la curva en dirección contraria a las agujas del reloj. \\
\end{defin}

\begin{rmk}La definición formal de orientación positiva es complicada y no forma parte de este curso.
\end{rmk}





\newpage

\section{Teorema de Green}

Ahora estamos listos para establecer el teorema de Green \\

\begin{thm}Sea $C$ una curva suave a trozos, cerrada y orientada positivamente y sea $F=(M,N)$ donde $M$y $N$ tienen derivadas parciales continuas en una región abierta que contiene a la región $R$ limitada por $C$, entonces
\end{thm}


\begin{align*}
    \oint_C F \cdot dr = \iint_R \left( N_y-M_x\right) dA
\end{align*}

Note que si $F$ es conservativo entonces necesariamente $ \displaystyle \oint_C F \cdot dr = 0$ \\




\begin{ejem}Empleando el Teorema de Green determine el área de la región limitada por la hipocicloide dada en forma paramétrica por $r(t) =(\cos^3(t),\sin^3(t))$, $0 \leq t \leq 2\pi$ \\
\end{ejem}

\begin{sol}Tenemos que el área será
\end{sol}

\begin{align*}
    \iint_R 1 dA
\end{align*}

Donde $R$ es la región limitada por la hipocicloide, para poder usar Green necesitamos una función $F=(M,N)$ que cumpla con que $N_x-M_y=1$, haciendo $F=(0,x)$, funciona, luego

\begin{align*}
    \iint_R 1 dA &=  \oint_C (0,x) \cdot r \\
    &=  \int_0^{2\pi} (0,\cos^3(t)) \cdot (-3\cos^2(t)\sin(t),3\sin^2(t)\cos(t)) dt\\
    &=  \int_0^{2\pi} 3\cos^4(t)\sin^2(t)dt\\
    &=  3\int_0^{2\pi} \cos^4(t)(1-\cos^2(t)))dt\\
    &=  3\int_0^{2\pi} \cos^4(t)-\cos^6(t)dt\\
    &= 3 \frac{\pi}{8} \\
    &= \frac{3\pi}{8} \\
\end{align*}

















\end{document}